\documentclass[12pt,a4paper]{article}
\usepackage[utf8]{inputenc}
\usepackage[T5]{fontenc}
\usepackage[vietnamese]{babel}
\usepackage{amsmath, amssymb}
\usepackage{array}
\usepackage{enumitem}
\usepackage{setspace}
\setstretch{1.08}
\usepackage{graphicx}
\usepackage{ifthen}
\usepackage{tikz}
\usetikzlibrary{arrows.meta, positioning, patterns}
\IfFileExists{pgf-pie.sty}{\usepackage{pgf-pie}}{}
\usepackage{pgfplots}
\pgfplotsset{compat=1.18}
\usepackage{chemfig}
\usepackage{longtable,booktabs}
\usepackage{adjustbox}
\providecommand{\errorpart}[2]{$\underset{\textbf{#1}}{\underline{\text{#2}}}$}
\usepackage[top=0.8in, left=0.75in, right=0.75in, bottom=0.85in, headheight=38pt, headsep=18pt, footskip=1.2cm]{geometry}
\usepackage{fancyhdr}
\usepackage[colorlinks=true, linkcolor=blue, urlcolor=blue]{hyperref}

\setlength{\parindent}{0pt}
\setlength{\parskip}{5pt plus 2pt minus 1pt}

\setlist[enumerate]{itemsep=4pt, parsep=2pt, topsep=4pt, partopsep=2pt}
\setlist[itemize]{itemsep=4pt, parsep=2pt, topsep=4pt}

\makeatletter
\AtEndDocument{\immediate\write\@auxout{\string\newlabel{LecturingLastPage}{{}{\thepage}{}{}{}}}}
\AtEndDocument{\immediate\write\@auxout{\string\newlabel{LecturingPageCount}{{}{\two@digits{\value{page}}}{}{}{}}}}
\makeatother
\pagestyle{fancy}
\fancyhf{}
\cfoot{\small Trang \thepage/\pageref*{LecturingLastPage} - Mã đề thi 2615}
\renewcommand{\headrulewidth}{0pt}
\renewcommand{\footrulewidth}{0pt}

\usepackage{tabularx}
\begin{document}

\sloppy

\noindent
\begin{minipage}[t]{0.37\textwidth}
    \small
    \centering
    \IfFileExists{logo.png}{\includegraphics[height=0.9cm]{logo.png}\\[0.1cm]}{%
      \IfFileExists{logo.jpg}{\includegraphics[height=0.9cm]{logo.jpg}\\[0.1cm]}{}%
    }
    {\bfseries \mbox{}}\\[-0.2cm]
    \rule{3cm}{0.5pt}\\[0.1cm]
    {\bfseries ĐỀ LUYỆN TẬP -- BẢN NGHIÊN CỨU}\\{\small (Đề thi có \pageref*{LecturingPageCount} trang)}
\end{minipage}%
\hfill
\begin{minipage}[t]{0.62\textwidth}
    \small
    \centering
    {\bfseries \sbox0{KỲ THI TỐT NGHIỆP TRUNG HỌC PHỔ THÔNG NĂM 2026}\ifdim\wd0>\linewidth KỲ THI TỐT NGHIỆP TRUNG HỌC PHỔ~THÔNG\\NĂM 2026\else\usebox0\fi}\\[0.15cm]
    {\bfseries Môn thi: HOÁ HỌC}\\[0.15cm]
    {\itshape Thời gian làm bài: 50 phút, không kể thời gian phát đề}
\end{minipage}

\vspace{0.4cm}
\noindent
\begin{minipage}[t]{0.68\textwidth}
    Họ, tên thí sinh: \dotfill\\[0.15cm]
    Số báo danh: \dotfill
\end{minipage}%
\hfill
\fbox{\begin{minipage}{3.6cm}\centering\textbf{Mã đề: 2615}\end{minipage}}

\vspace{0.5cm}

% LECTURING-PART: P1
\textbf{Phần I: Trắc nghiệm nhiều lựa chọn} \textit{Thí sinh trả lời từ câu 1 đến câu 18. Mỗi câu hỏi thí sinh chỉ chọn một phương án.}


\par\noindent\begin{minipage}{\linewidth}
\noindent \textbf{Câu 1:} Kim loại nào sau đây tác dụng mạnh với nước ở nhiệt độ thường, tạo thành dung dịch base và khí hydrogen?\\
\noindent\makebox[0.24\textwidth][l]{A. $\text{Ag}$.}\makebox[0.24\textwidth][l]{B. $\text{Cu}$.}\makebox[0.24\textwidth][l]{C. $\text{Na}$.}\makebox[0.24\textwidth][l]{D. $\text{Fe}$.}
\end{minipage}\par\vspace{0.12cm}

\par\noindent\begin{minipage}{\linewidth}
\noindent \textbf{Câu 2:} Chất nào sau đây là amine bậc hai?\\
\noindent A. $\text{CH}_3\text{NH}_2$. \hfill B. $(\text{CH}_3)_2\text{NH}$. \hfill C. $(\text{CH}_3)_3\text{N}$. \hfill D. $\text{C}_6\text{H}_5\text{NH}_2$.\\
\end{minipage}\par\vspace{0.12cm}

\par\noindent\begin{minipage}{\linewidth}
\noindent \textbf{Câu 3:} Công thức Haworth của monosaccharide X được biểu diễn ở hình bên dưới.\\
\begin{center}
\begin{adjustbox}{max width=0.95\linewidth, max totalheight=0.38\textheight}
\begin{tikzpicture}[scale=0.9, every node/.style={font=\small}]
\coordinate (C2) at (2,0); \coordinate (C3) at (1.2,-1);
\coordinate (C4) at (-0.4,-1); \coordinate (C5) at (-1.2,0);
\node (O) at (0.4,1) {$\mathrm{O}$};
\draw[thick] (O) -- (C2) -- (C3) -- (C4) -- (C5) -- (O);
\draw[line width=1.8pt] (C3)--(C4);
\draw (C2)--(2,1.1) node[above] {$\mathrm{OH}$};
\draw (C2)--(2,-1.8) node[below right,inner sep=1pt] {$\mathrm{CH_2OH}$};
\draw (C3)--(1.2,-0.05) node[above] {$\mathrm{OH}$};
\draw (C3)--(1.2,-1.9) node[below] {$\mathrm{H}$};
\draw (C4)--(-0.4,-1.9) node[below] {$\mathrm{OH}$};
\draw (C4)--(-0.4,-0.05) node[above] {$\mathrm{H}$};
\draw (C5)--(-1.2,1.1) node[above] {$\mathrm{CH_2OH}$};
\draw (C5)--(-1.2,-0.9) node[below] {$\mathrm{H}$};
\end{tikzpicture}
\end{adjustbox}
\end{center}
Chất X là\\
\noindent\makebox[0.24\textwidth][l]{A. fructose.}\makebox[0.24\textwidth][l]{B. maltose.}\makebox[0.24\textwidth][l]{C. sucrose.}\makebox[0.24\textwidth][l]{D. glucose.}
\end{minipage}\par\vspace{0.12cm}

\par\noindent\begin{minipage}{\linewidth}
\noindent \textbf{Câu 4:} Cho bảng giá trị thế điện cực chuẩn của một số cặp oxi hoá -- khử ở $25\,^{\circ}\text{C}$ như sau:\\
\begin{center}
\begin{adjustbox}{max width=\linewidth}
\begin{tabular}{|c|c|c|c|c|}
\hline
Cặp oxi hoá -- khử & $\text{Mg}^{2+}/\text{Mg}$ & $\text{Zn}^{2+}/\text{Zn}$ & $\text{Cu}^{2+}/\text{Cu}$ & $\text{Ag}^{+}/\text{Ag}$ \\
\hline
$E^{\circ}$ (V) & $-2{,}37$ & $-0{,}76$ & $+0{,}34$ & $+0{,}80$ \\
\hline
\end{tabular}
\end{adjustbox}
\end{center}
Sức điện động chuẩn của pin điện hoá $\text{Zn}$ -- $\text{Ag}$ là\\
\noindent\makebox[0.24\textwidth][l]{A. $0{,}04$ V.}\makebox[0.24\textwidth][l]{B. $1{,}10$ V.}\makebox[0.24\textwidth][l]{C. $1{,}56$ V.}\makebox[0.24\textwidth][l]{D. $2{,}36$ V.}
\end{minipage}\par\vspace{0.12cm}

\par\noindent\begin{minipage}{\linewidth}
\noindent \textbf{Câu 5:} Thuỷ phân hoàn toàn ethyl acetate trong dung dịch $\text{NaOH}$ đun nóng, thu được các sản phẩm hữu cơ là\\
\noindent A. $\text{CH}_3\text{COONa}$ và $\text{C}_2\text{H}_5\text{OH}$. \hfill B. $\text{C}_2\text{H}_5\text{COONa}$ và $\text{CH}_3\text{OH}$.\\
\noindent C. $\text{HCOONa}$ và $\text{C}_3\text{H}_7\text{OH}$. \hfill D. $\text{CH}_3\text{COOH}$ và $\text{C}_2\text{H}_5\text{ONa}$.\\
\end{minipage}\par\vspace{0.12cm}

\par\noindent\begin{minipage}{\linewidth}
\noindent\textbf{Câu 6:} Trong bình kín nhiệt độ không đổi tồn tại cân bằng N$_{2}$O$_{4}$(g) \(\rightleftharpoons\) 2NO$_{2}$(g). Nếu giảm thể tích bình đột ngột và đợi cân bằng mới, xu hướng dịch chuyển ban đầu là gì?\par
\noindent A. Về phải vì tăng số mol khí.\par
\noindent B. Về trái vì giảm số mol khí.\par
\noindent C. Không đổi vì nhiệt độ cố định.\par
\noindent D. Không xác định vì thiếu xúc tác.\par
\end{minipage}\par\vspace{0.12cm}

\par\noindent\begin{minipage}{\linewidth}
\noindent \textbf{Câu 7:} Polymer nào sau đây được điều chế bằng phản ứng trùng ngưng?\\
\noindent\makebox[0.48\textwidth][l]{A. Polystyrene.}\makebox[0.48\textwidth][l]{B. Poly(vinyl chloride).}\\
\noindent\makebox[0.48\textwidth][l]{C. Nylon-6,6.}\makebox[0.48\textwidth][l]{D. Polyethylene.}
\end{minipage}\par\vspace{0.12cm}

\par\noindent\begin{minipage}{\linewidth}
\noindent \textbf{Câu 8:} Dung dịch của chất nào sau đây làm quỳ tím chuyển sang màu xanh?\\
\noindent\makebox[0.24\textwidth][l]{A. Lysine.}\makebox[0.24\textwidth][l]{B. Glutamic acid.}\makebox[0.24\textwidth][l]{C. Glycine.}\makebox[0.24\textwidth][l]{D. Alanine.}
\end{minipage}\par\vspace{0.12cm}

\par\noindent\begin{minipage}{\linewidth}
\noindent \textbf{Câu 9:} Nhỏ từ từ đến dư dung dịch $\text{NH}_3$ vào ống nghiệm đựng dung dịch $\text{CuSO}_4$. Hiện tượng quan sát được là\\
\noindent A. xuất hiện kết tủa màu xanh, sau đó kết tủa tan tạo dung dịch màu xanh lam đậm.\\
\noindent B. xuất hiện kết tủa màu xanh, kết tủa không tan khi thêm $\text{NH}_3$ dư.\\
\noindent C. xuất hiện kết tủa màu trắng, sau đó kết tủa tan tạo dung dịch không màu.\\
\noindent D. có khí mùi khai thoát ra và dung dịch chuyển sang không màu.
\end{minipage}\par\vspace{0.12cm}

\par\noindent\begin{minipage}{\linewidth}
\noindent \textbf{Câu 10:} Nguyên tử của nguyên tố X có số hiệu nguyên tử $Z = 17$. Vị trí của X trong bảng tuần hoàn là\\
\noindent\makebox[0.48\textwidth][l]{A. chu kì 3, nhóm VA.}\makebox[0.48\textwidth][l]{B. chu kì 4, nhóm IA.}\\
\noindent\makebox[0.48\textwidth][l]{C. chu kì 2, nhóm VIIA.}\makebox[0.48\textwidth][l]{D. chu kì 3, nhóm VIIA.}
\end{minipage}\par\vspace{0.12cm}

\par\noindent\begin{minipage}{\linewidth}
\noindent \textbf{Câu 11:} Điện phân dung dịch $\text{CuSO}_4$ với các điện cực trơ. Quá trình xảy ra ở anode là\\
\noindent A. sự oxi hoá phân tử $\text{H}_2\text{O}$.\\
\noindent B. sự oxi hoá ion $\text{SO}_4^{2-}$.\\
\noindent C. sự khử ion $\text{Cu}^{2+}$.\\
\noindent D. sự khử phân tử $\text{H}_2\text{O}$.
\end{minipage}\par\vspace{0.12cm}

\par\noindent\begin{minipage}{\linewidth}
\noindent \textbf{Câu 12:} Một mẫu nước giếng chứa các ion $\text{Ca}^{2+}$, $\text{Mg}^{2+}$, $\text{Cl}^{-}$ và $\text{SO}_4^{2-}$. Hoá chất nào sau đây làm mềm được mẫu nước này?\\
\noindent\makebox[0.48\textwidth][l]{A. $\text{Ca(OH)}_2$.}\makebox[0.48\textwidth][l]{B. $\text{Na}_2\text{CO}_3$.}\\
\noindent\makebox[0.48\textwidth][l]{C. $\text{HCl}$.}\makebox[0.48\textwidth][l]{D. $\text{NaCl}$.}
\end{minipage}\par\vspace{0.12cm}

\par\noindent\begin{minipage}{\linewidth}
\noindent \textbf{Câu 13:} Một ống thép (thành phần chính là $\text{Fe}$) được nối với một thanh kim loại M rồi cùng đặt trong đất ẩm. Với kim loại M nào sau đây thì ống thép bị ăn mòn điện hoá nhanh hơn?\\
\begin{center}
\begin{adjustbox}{max width=0.8\linewidth, max totalheight=0.25\textheight}
\begin{tikzpicture}[scale=0.9]
  \fill[brown!25] (0,0) rectangle (9,3);
  \draw[thick] (0,0) rectangle (9,3);
  \node[fill=brown!25, inner sep=2pt] at (4.5,0.45) {đất ẩm};
  \fill[gray!55] (0.8,1.2) rectangle (4.2,1.9);
  \draw (0.8,1.2) rectangle (4.2,1.9);
  \node[above] at (2.5,1.95) {ống thép (Fe)};
  \fill[gray!30] (6.4,1.1) rectangle (7.6,2.0);
  \draw (6.4,1.1) rectangle (7.6,2.0);
  \node[above] at (7.0,2.05) {kim loại M};
  \draw[thick] (4.2,1.55) -- (6.4,1.55);
  \node[below, font=\small] at (5.3,1.5) {dây dẫn};
\end{tikzpicture}
\end{adjustbox}
\end{center}
\noindent\makebox[0.24\textwidth][l]{A. $\text{Cu}$.}\makebox[0.24\textwidth][l]{B. $\text{Mg}$.}\makebox[0.24\textwidth][l]{C. $\text{Zn}$.}\makebox[0.24\textwidth][l]{D. $\text{Al}$.}
\end{minipage}\par\vspace{0.12cm}

\par\noindent\begin{minipage}{\linewidth}
\noindent\textbf{Câu 14:} Cho \(\Delta\)fH\(^{\circ}\)(CO$_{2}$(g))=-393,5 kJ/mol và \(\Delta\)fH\(^{\circ}\)(CO(g))=-110,5 kJ/mol; \(\Delta\)fH\(^{\circ}\)(O$_{2}$(g))=0. \(\Delta\)rH\(^{\circ}\) của CO(g)+½O$_{2}$(g)\(\rightarrow\)CO$_{2}$(g) là bao nhiêu?\par
\noindent A. -504,0 kJ/mol\par
\noindent B. +283,0 kJ/mol\par
\noindent C. -283,0 kJ/mol\par
\noindent D. +504,0 kJ/mol\par
\end{minipage}\par\vspace{0.12cm}

\par\noindent\begin{minipage}{\linewidth}
\noindent\textbf{Câu 15:} Cho propene phản ứng với HBr trong điều kiện không có peroxide. Sản phẩm hữu cơ chính theo quy tắc Markovnikov là chất nào?\par
\noindent A. 1-bromopropane\par
\noindent B. 2-bromopropane\par
\noindent C. 1,2-dibromopropane\par
\noindent D. Propane\par
\end{minipage}\par\vspace{0.12cm}

\par\noindent\begin{minipage}{\linewidth}
\noindent\textbf{Câu 16:} Phổ hồng ngoại của chất hữu cơ lỏng X có vạch hấp thụ mạnh gần 1715 cm$^{-1}$; không có dải rộng O--H vùng 3200--3600 cm$^{-1}$. Chất nào phù hợp nhất trong các lựa chọn dưới đây?\par
\noindent A. Propan-2-ol\par
\noindent B. Propanone\par
\noindent C. Propan-1-amine\par
\noindent D. Propane\par
\end{minipage}\par\vspace{0.12cm}

\par\noindent\begin{minipage}{\linewidth}
\noindent \textbf{Câu 17:} Phát biểu nào sau đây về protein là \textbf{đúng}?\\
\noindent A. Protein đơn giản khi thuỷ phân hoàn toàn chỉ thu được các $\alpha$-amino acid.\\
\noindent B. Protein không bị đông tụ khi đun nóng.\\
\noindent C. Liên kết peptide là liên kết giữa nhóm $\text{-COOH}$ và nhóm $\text{-OH}$.\\
\noindent D. Tất cả các protein đều tan tốt trong nước tạo dung dịch keo.
\end{minipage}\par\vspace{0.12cm}

\par\noindent\begin{minipage}{\linewidth}
\noindent \textbf{Câu 18:} Hoà tan hoàn toàn $7{,}8$ gam hỗn hợp gồm $\text{Mg}$ và $\text{Al}$ trong dung dịch $\text{HCl}$ dư, thu được $9{,}916$ lít khí $\text{H}_2$ ở điều kiện chuẩn (một mol khí chiếm $24{,}79$ lít). Cho nguyên tử khối: $\text{Mg} = 24$; $\text{Al} = 27$. Khối lượng $\text{Al}$ trong hỗn hợp ban đầu là\\
\noindent\makebox[0.24\textwidth][l]{A. $2{,}4$ gam.}\makebox[0.24\textwidth][l]{B. $2{,}7$ gam.}\makebox[0.24\textwidth][l]{C. $4{,}05$ gam.}\makebox[0.24\textwidth][l]{D. $5{,}4$ gam.}

\vspace{0.5cm}
\end{minipage}\par\vspace{0.12cm}
% LECTURING-PART: P2
\textbf{Phần II: Trắc nghiệm đúng/sai} \textit{Thí sinh trả lời từ câu 1 đến câu 4. Trong mỗi ý a), b), c), d) ở mỗi câu, thí sinh chọn đúng hoặc sai.}


\par\noindent\begin{minipage}{\linewidth}
\noindent \textbf{Câu 1:} Một nhóm học sinh tiến hành điện phân dung dịch $\text{CuSO}_4$ với hai điện cực graphite theo sơ đồ ở hình bên dưới, cường độ dòng điện không đổi $I = 1{,}93$ A trong thời gian $2000$ giây. Biết hiệu suất điện phân là $100\%$, hằng số Faraday $F = 96485$ C/mol, nguyên tử khối $\text{Cu} = 64$, và trong suốt thời gian điện phân ion $\text{Cu}^{2+}$ vẫn còn dư.\\
\begin{center}
\begin{adjustbox}{max width=0.95\linewidth, max totalheight=0.21\textheight}
\begin{tikzpicture}[scale=0.9]
  % cốc
  \draw[thick] (0,4) -- (0,0) -- (6,0) -- (6,4);
  % mặt dung dịch
  \fill[gray!15] (0.02,0.02) rectangle (5.98,3);
  \draw (0,3) -- (6,3);
  % điện cực
  \fill[gray!70] (1.3,1.0) rectangle (1.7,4.6);
  \draw (1.3,1.0) rectangle (1.7,4.6);
  \fill[gray!70] (4.3,1.0) rectangle (4.7,4.6);
  \draw (4.3,1.0) rectangle (4.7,4.6);
  % dây dẫn và nguồn
  \draw (1.5,4.6) -- (1.5,5.8) -- (2.85,5.8);
  \draw (4.5,4.6) -- (4.5,5.8) -- (3.15,5.8);
  \draw[very thick] (2.85,5.55) -- (2.85,6.05); % vạch ngắn: cực âm
  \draw[very thick] (3.15,5.35) -- (3.15,6.25); % vạch dài: cực dương


  % nhãn

  \draw (-0.1,3.8) -- (1.3,3.8);

  \draw (6.1,3.8) -- (4.7,3.8);
    % captions last so nothing paints over them
    % captions last so nothing paints over them
    \node[above] at (2.6,6.0) {$-$};
\node[above] at (3.4,6.2) {$+$};
\node[left] at (-0.1,3.8) {cathode};
\node[right] at (6.1,3.8) {anode};
\node[fill=gray!15, inner sep=2pt] at (3,0.5) {dung dịch $\text{CuSO}_4$};
\end{tikzpicture}
\end{adjustbox}
\end{center}
\noindent a) Trong quá trình điện phân, ở điện cực nối với cực âm của nguồn có lớp kim loại màu đỏ bám vào.\\

\noindent b) Khí thoát ra ở anode là khí hydrogen.\\

\noindent c) Khối lượng kim loại bám vào cathode sau $2000$ giây xấp xỉ $1{,}28$ gam.\\

\noindent d) Giá trị pH của dung dịch tăng dần trong quá trình điện phân.\\
\end{minipage}\par\vspace{0.12cm}

\par\noindent\begin{minipage}{\linewidth}
\noindent \textbf{Câu 2:} Tiến hành thí nghiệm phản ứng tráng bạc của glucose theo các bước sau:\\
Bước 1: Cho $2$ mL dung dịch $\text{AgNO}_3$ $1\%$ vào ống nghiệm sạch, nhỏ từ từ dung dịch $\text{NH}_3$ $5\%$ vào, lắc đều đến khi kết tủa vừa tạo thành tan hết.\\
Bước 2: Thêm tiếp $1$ mL dung dịch glucose $10\%$ vào ống nghiệm.\\
Bước 3: Đặt ống nghiệm vào cốc nước nóng khoảng $60$ -- $70\,^{\circ}\text{C}$ trong vài phút.\\
Cho nguyên tử khối: $\text{H} = 1$; $\text{C} = 12$; $\text{O} = 16$; $\text{Ag} = 108$.\\
\noindent a) Sau bước 3, trên thành ống nghiệm xuất hiện một lớp kim loại sáng như gương.\\

\noindent b) Trong phản ứng ở bước 3, glucose đóng vai trò chất oxi hoá.\\

\noindent c) Nếu thay dung dịch glucose ở bước 2 bằng dung dịch sucrose thì sau bước 3 cũng thu được lớp kim loại sáng như gương.\\

\noindent d) Trong một thí nghiệm khác, cho $0{,}9$ gam glucose phản ứng hoàn toàn với dung dịch $\mathrm{AgNO_3/NH_3}$ dư thì thu được $1{,}08$ gam Ag.\\
\end{minipage}\par\vspace{0.12cm}

\par\noindent\begin{minipage}{\linewidth}
\noindent \textbf{Câu 3:} Tiến hành thí nghiệm xà phòng hoá chất béo theo các bước sau:\\
Bước 1: Cho khoảng $2$ gam mỡ lợn và $5$ mL dung dịch $\text{NaOH}$ $40\%$ vào bát sứ.\\
Bước 2: Đun sôi nhẹ hỗn hợp, liên tục khuấy đều bằng đũa thuỷ tinh, thỉnh thoảng thêm vài giọt nước cất.\\
Bước 3: Sau khoảng $8$ -- $10$ phút, rót thêm vào hỗn hợp $10$ mL dung dịch $\text{NaCl}$ bão hoà nóng, khuấy nhẹ rồi để nguội.\\
Cho nguyên tử khối: $\text{H} = 1$; $\text{C} = 12$; $\text{O} = 16$; $\text{Na} = 23$.\\
\noindent a) Sau bước 3, có một lớp chất rắn màu trắng nổi lên trên bề mặt hỗn hợp.\\

\noindent b) Mục đích chính của việc thêm dung dịch $\text{NaCl}$ bão hoà ở bước 3 là làm tăng tốc độ phản ứng xà phòng hoá.\\

\noindent c) Việc thêm nước cất ở bước 2 giúp hỗn hợp không bị cạn khô, nhờ đó phản ứng xà phòng hoá tiếp tục xảy ra.\\

\noindent d) Xà phòng hoá hoàn toàn $8{,}9$ gam tristearin bằng dung dịch $\text{NaOH}$ dư thu được $2{,}76$ gam glycerol.\\
\end{minipage}\par\vspace{0.12cm}

\par\noindent\begin{minipage}{\linewidth}
\noindent \textbf{Câu 4:} Để xác định nồng độ của một dung dịch $\text{NaOH}$, một học sinh lấy chính xác $10{,}0$ mL dung dịch này cho vào bình tam giác, thêm 2 giọt phenolphthalein, rồi chuẩn độ bằng dung dịch $\text{HCl}$ $0{,}100$ M. Kết quả ba lần chuẩn độ được ghi ở bảng sau:\\
\begin{center}
\begin{adjustbox}{max width=\linewidth}
\begin{tabular}{|c|c|c|c|}
\hline
Lần chuẩn độ & 1 & 2 & 3 \\
\hline
Thể tích dung dịch $\text{HCl}$ đã dùng (mL) & $10{,}1$ & $9{,}9$ & $10{,}0$ \\
\hline
\end{tabular}
\end{adjustbox}
\end{center}
\noindent a) Thời điểm kết thúc chuẩn độ là khi dung dịch trong bình tam giác chuyển từ màu hồng sang không màu và giữ nguyên trong khoảng 20 giây.\\

\noindent b) Thể tích trung bình của dung dịch $\text{HCl}$ đã dùng trong ba lần chuẩn độ là $10{,}0$ mL.\\

\noindent c) Nồng độ của dung dịch $\text{NaOH}$ ban đầu là $0{,}010$ M.\\

\noindent d) Dung dịch $\text{HCl}$ trong thí nghiệm này được chứa trong pipette và nhỏ từ từ vào bình tam giác.\\

\vspace{0.5cm}
\end{minipage}\par\vspace{0.12cm}
% LECTURING-PART: P3
\textbf{Phần III: Trắc nghiệm trả lời ngắn} \textit{Thí sinh trả lời từ câu 1 đến câu 6. Thí sinh chỉ điền một con số (có thể có dấu âm, dấu phẩy thập phân) vào ô trống.}


\par\noindent\begin{minipage}{\linewidth}
\noindent \textbf{Câu 1:} Cho các polymer sau: (1) polyethylene; (2) nylon-6,6; (3) cellulose; (4) poly(vinyl chloride); (5) polyacrylonitrile. Có bao nhiêu polymer trong danh sách được điều chế bằng phản ứng trùng hợp?\\
\end{minipage}\par\vspace{0.12cm}

\par\noindent\begin{minipage}{\linewidth}
\noindent \textbf{Câu 2:} Có bao nhiêu ester mạch hở có cùng công thức phân tử $\text{C}_4\text{H}_8\text{O}_2$?\\
\end{minipage}\par\vspace{0.12cm}

\par\noindent\begin{minipage}{\linewidth}
\noindent\textbf{Câu 3:} Ngữ liệu và mô hình công nghệ tự biên soạn. Một nhóm học sinh thiết kế một mẻ điều chế hỗn hợp sodium oleate và sodium stearate. Nhóm được chọn lượng nguyên liệu sử dụng trong giới hạn sau:\par
\begin{center}\small\begin{tabularx}{\linewidth}{|X|X|}
\hline
Nguyên liệu & Lượng tối đa được sử dụng \\ \hline
Triolein & 0,160 mol \\ \hline
Tristearin & 0,030 mol \\ \hline
Hydrogen & 0,090 mol \\ \hline
Sodium hydroxide & 0,540 mol \\ \hline
\end{tabularx}\end{center}
Trước khi xà phòng hóa, nhóm có thể hydrogen hóa một phần bất kì số liên kết C=C của triolein đã chọn. Sau đó xà phòng hóa hoàn toàn lượng triglyceride đã chọn bằng sodium hydroxide. Mỗi liên kết C=C được hydrogen hóa tiêu thụ 1 mol hydrogen trên 1 mol liên kết. Mỗi mol triglyceride tiêu thụ 3 mol sodium hydroxide và tạo 3 mol muối carboxylate. Hydrogen hóa biến gốc oleate thành gốc stearate; các liên kết ester không bị ảnh hưởng.\par
Hỗn hợp muối cuối cùng phải có số mol sodium stearate chiếm từ 35,0\% đến 45,0\% tổng số mol hai muối. Nhóm muốn thu được khối lượng sodium oleate lớn nhất có thể, không nhất thiết sử dụng hết từng nguyên liệu.\par
Giả sử mọi phản ứng đạt hiệu suất 100\%, không có phản ứng phụ hoặc hao hụt; nguyên liệu tinh khiết và có thể chia lượng liên tục. Chỉ tính khối lượng muối khan, không tính glycerol, nước hoặc chất dư. Cho M(sodium oleate) = 304 g/mol.\par
Khối lượng sodium oleate lớn nhất thu được là bao nhiêu gam? Làm tròn đến 0,1 g.\par
\end{minipage}\par\vspace{0.12cm}

\par\noindent\begin{minipage}{\linewidth}
\noindent\textbf{Câu 4:} Tình huống và số liệu tự biên soạn. Có 10 phần dung dịch giống nhau; mỗi phần chứa 0,010 mol glucose và 0,005 mol sucrose. Mỗi phần phải được xử lí nguyên phần theo một trong hai quy trình, hoặc không sử dụng:\par
\begin{center}\small\begin{tabularx}{\linewidth}{|X|X|X|X|}
\hline
Quy trình & Xử lí carbohydrate & Lượng AgNO3 cấp cố định cho mỗi phần & Hiệu suất thu hồi silver \\ \hline
R & Không thủy phân & 0,020 mol & 80,0\% \\ \hline
H & Thủy phân hoàn toàn sucrose, sau đó trung hòa & 0,040 mol & 76,0\% \\ \hline
\end{tabularx}\end{center}
Thuốc thử chứa AgNO3 trong môi trường ammonia được chuẩn bị phù hợp. Trong điều kiện này, 1 mol glucose hoặc fructose có thể tạo tối đa 2 mol silver; sucrose chưa thủy phân không tạo silver. Hiệu suất trong bảng tính theo lượng silver tối đa từ carbohydrate của phần dung dịch tương ứng. Toàn bộ AgNO3 đã cấp cho một phần được tính là đã sử dụng, kể cả phần không được thu hồi thành silver; không tái sử dụng thuốc thử.\par
Tổng lượng AgNO3 có thể cấp là 0,260 mol. Silver thu hồi từ các phần được gộp lại để tạo lớp phủ, giả sử công đoạn tạo lớp phủ không hao hụt. Mỗi tấm cần phủ trên tổng diện tích 180 cm$^{2}$ một lớp silver dày đều 1,00 \(\mu\)m. Bỏ qua thể tích ở mép và độ nhám; khối lượng riêng của silver là 10,5 g/cm$^{3}$. Cho M(Ag) = 108 g/mol và 1 \(\mu\)m = 10$^{-4}$ cm.\par
Có thể hoàn thành nhiều nhất bao nhiêu tấm? Chỉ tính tấm được phủ đủ; ghi một số nguyên, đơn vị tấm.\par
\end{minipage}\par\vspace{0.12cm}

\par\noindent\begin{minipage}{\linewidth}
\noindent\textbf{Câu 5:} Tình huống tự biên soạn. Một chất X có công thức phân tử C9H13N, chứa đúng một vòng benzene và không chứa vòng nào khác. Cần xác định tất cả công thức cấu tạo của X phù hợp đồng thời với các dữ kiện sau:\par
\(\bullet\) X là amine bậc một; nguyên tử nitrogen liên kết trực tiếp với một nguyên tử carbon no thuộc mạch nhánh, không liên kết trực tiếp với vòng benzene.\par
\(\bullet\) Khi oxi hóa hoàn toàn các mạch nhánh bằng potassium permanganate dư, nóng, rồi acid hóa, sản phẩm hữu cơ chứa vòng benzene duy nhất là benzene-1,4-dicarboxylic acid.\par
\(\bullet\) Để xét phép oxi hóa này, sử dụng mô hình: mỗi mạch nhánh gắn với vòng qua một carbon có ít nhất một hydrogen được chuyển thành một nhóm --COOH tại chính vị trí gắn đó; các carbon xa vòng và nitrogen không nằm trong sản phẩm hữu cơ chứa vòng. Vị trí tương đối của các nhóm thế trên vòng không thay đổi.\par
Có bao nhiêu công thức cấu tạo của X phù hợp? Chỉ đếm đồng phân cấu tạo, không đếm riêng các đồng phân quang học; ghi một số nguyên.\par
\end{minipage}\par\vspace{0.12cm}

\par\noindent\begin{minipage}{\linewidth}
\noindent\textbf{Câu 6:} Một mô hình trùng ngưng mạch thẳng chỉ tạo liên kết amide, không tạo vòng và không có phản ứng phụ. Cấp 0,120 mol H2N--R--NH2, 0,100 mol HOOC--R\(\prime\)--COOH và 0,0300 mol CH3COOH (chất chặn đầu mạch một chức). Tất cả nhóm --COOH phản ứng hết, --NH2 được phép dư. Mỗi liên kết amide hình thành giữa hai phân tử đang tách rời và làm mất một H2O. Mỗi phân tử hữu cơ cuối là một mạch thẳng có hai đầu; đầu mạch chỉ có thể là --NH2 tự do hoặc đầu acetyl từ CH3COOH. Monomer còn dư cũng được tính là một phân tử mạch thẳng. Không áp đặt chiều dài tối thiểu; các cách phân bố liên kết phù hợp bảo toàn đều được phép. Số mol phân tử có CẢ HAI đầu là acetyl ít nhất bằng bao nhiêu? Ghi mol, 3 chữ số thập phân.\par
\end{minipage}\par\vspace{0.12cm}
\begin{center}
{\bfseries -{}-{}-{}-{}-{}-}\ \textbf{HẾT}\ {\bfseries -{}-{}-{}-{}-{}-}
\end{center}
\vspace{0.15cm}
{\small\itshape
Thí sinh không được sử dụng tài liệu. Cán bộ coi thi không giải thích gì thêm.}

\end{document}
